\clearpage
\section{Scriptural Date and Location}\label{sec:date-location}
% Generated. Do not edit.
% Deterministic projection of research/chronology.toml.
% schema: 3
% document: liturgy/roman-rite/1962/propers/temporal/59-nineteenth-after-pentecost
% generated-by: tools/tpt proper-chronology annotations
\newcommand{\chronologyannotationclaim}[6]{#6}
\newcommand{\chronologyannotationcomparisonclaim}[8]{#8}
\newcommand{\chronologyannotationreach}[3]{}
\newcommand{\chronologyannotationgroup}[3]{#3}
\newcommand{\chronologyannotationcomparisongroup}[4]{#4}
\newcommand{\chronologyannotation}[1]{%
  \ifcsname triptychchronologyannotation@#1\endcsname
    \csname triptychchronologyannotation@#1\endcsname
  \else
    \PackageError{triptych}{No chronology annotation for #1}{}%
  \fi}
% comparison-dependency: src/gpt/liturgy/roman-rite/1962/propers/temporal/59-nineteenth-after-pentecost/research/chronology-profile-comparisons.toml sha256=7306dbbfa42399c1dab06cac21ca7fbdbf67fd0cb1fef34297adddca32b80cbc
% introit: Introit
\expandafter\def\csname triptychchronologyannotation@introit\endcsname{%
  \chronologyannotationgroup{composition}{preferred}{\textbf{Composition}: \chronologyannotationreach{Ps.77.1}{Ps}{true}\chronologyannotationclaim{critical.psalms.latest-composition-boundary}{composition}{catholic-critical-v1}{preferred}{before the Maccabean period, around 165 B.C.; no individual psalm can be dated securely}{Before c. 165 B.C.}}%
}
% epistle: Epistle
\expandafter\def\csname triptychchronologyannotation@epistle\endcsname{%
  \chronologyannotationgroup{composition}{disputed}{\textbf{Composition} -- disputed: \chronologyannotationreach{Eph.4.23}{Eph}{true}\chronologyannotationreach{Eph.4.24}{Eph}{true}\chronologyannotationreach{Eph.4.25}{Eph}{true}\chronologyannotationreach{Eph.4.26}{Eph}{true}\chronologyannotationreach{Eph.4.27}{Eph}{true}\chronologyannotationreach{Eph.4.28}{Eph}{true}\chronologyannotationclaim{composition.epistle-to-the-ephesians}{composition}{catholic-traditional-v1}{disputed}{a period between 58 and 63}{A.D. 58--63}; \chronologyannotationreach{Eph.4.23}{Eph}{true}\chronologyannotationreach{Eph.4.24}{Eph}{true}\chronologyannotationreach{Eph.4.25}{Eph}{true}\chronologyannotationreach{Eph.4.26}{Eph}{true}\chronologyannotationreach{Eph.4.27}{Eph}{true}\chronologyannotationreach{Eph.4.28}{Eph}{true}\chronologyannotationclaim{composition.epistle-to-the-ephesians}{composition}{catholic-traditional-v1}{disputed}{(Philemon; Colossians; Ephesians; Philippians), 61}{A.D. 61}.}%
  \space \chronologyannotationcomparisongroup{catholic-critical-v1}{composition}{composition-only}{\textbf{Comparison (catholic-critical-v1) -- Composition} -- composition-only: \chronologyannotationreach{Eph.4.23}{Eph}{true}\chronologyannotationreach{Eph.4.24}{Eph}{true}\chronologyannotationreach{Eph.4.25}{Eph}{true}\chronologyannotationreach{Eph.4.26}{Eph}{true}\chronologyannotationreach{Eph.4.27}{Eph}{true}\chronologyannotationreach{Eph.4.28}{Eph}{true}\chronologyannotationcomparisonclaim{critical.ephesians-later-disciple}{composition}{catholic-critical-v1}{catholic-critical-v1}{disputed}{passage.united-states-conference-of-catholic-bishops.new-american-bible-revised-edition.english-usccb-web-2026-09-28.ephesians-introduction}{around A.D. 80--100}{Ephesians under the later-disciple hypothesis (approximate), A.D. 80--100}.}%
}
% gradual: Gradual
\expandafter\def\csname triptychchronologyannotation@gradual\endcsname{%
  \chronologyannotationgroup{composition}{preferred}{\textbf{Composition}: \chronologyannotationreach{Ps.140.2}{Ps}{true}\chronologyannotationclaim{critical.psalms.latest-composition-boundary}{composition}{catholic-critical-v1}{preferred}{before the Maccabean period, around 165 B.C.; no individual psalm can be dated securely}{Before c. 165 B.C.}}%
}
% alleluia: Alleluia
\expandafter\def\csname triptychchronologyannotation@alleluia\endcsname{%
  \chronologyannotationgroup{composition}{preferred}{\textbf{Composition}: \chronologyannotationreach{Ps.104.1}{Ps}{true}\chronologyannotationclaim{critical.psalms.latest-composition-boundary}{composition}{catholic-critical-v1}{preferred}{before the Maccabean period, around 165 B.C.; no individual psalm can be dated securely}{Before c. 165 B.C.}}%
}
% gospel: Gospel
\expandafter\def\csname triptychchronologyannotation@gospel\endcsname{%
  \chronologyannotationgroup{narrated-event}{research-pending}{\textbf{Event} -- Narrated event date unresolved.}%
  \space \chronologyannotationgroup{composition}{disputed}{\textbf{Composition} -- disputed: \chronologyannotationreach{Matt.22.1}{Matt}{true}\chronologyannotationreach{Matt.22.10}{Matt}{true}\chronologyannotationreach{Matt.22.11}{Matt}{true}\chronologyannotationreach{Matt.22.12}{Matt}{true}\chronologyannotationreach{Matt.22.13}{Matt}{true}\chronologyannotationreach{Matt.22.14}{Matt}{true}\chronologyannotationreach{Matt.22.2}{Matt}{true}\chronologyannotationreach{Matt.22.3}{Matt}{true}\chronologyannotationreach{Matt.22.4}{Matt}{true}\chronologyannotationreach{Matt.22.5}{Matt}{true}\chronologyannotationreach{Matt.22.6}{Matt}{true}\chronologyannotationreach{Matt.22.7}{Matt}{true}\chronologyannotationreach{Matt.22.8}{Matt}{true}\chronologyannotationreach{Matt.22.9}{Matt}{true}\chronologyannotationclaim{composition.gospel-of-matthew}{composition}{catholic-traditional-v1}{disputed}{about A.D. 38-45}{c. A.D. 38--45}; \chronologyannotationreach{Matt.22.1}{Matt}{true}\chronologyannotationreach{Matt.22.10}{Matt}{true}\chronologyannotationreach{Matt.22.11}{Matt}{true}\chronologyannotationreach{Matt.22.12}{Matt}{true}\chronologyannotationreach{Matt.22.13}{Matt}{true}\chronologyannotationreach{Matt.22.14}{Matt}{true}\chronologyannotationreach{Matt.22.2}{Matt}{true}\chronologyannotationreach{Matt.22.3}{Matt}{true}\chronologyannotationreach{Matt.22.4}{Matt}{true}\chronologyannotationreach{Matt.22.5}{Matt}{true}\chronologyannotationreach{Matt.22.6}{Matt}{true}\chronologyannotationreach{Matt.22.7}{Matt}{true}\chronologyannotationreach{Matt.22.8}{Matt}{true}\chronologyannotationreach{Matt.22.9}{Matt}{true}\chronologyannotationclaim{composition.gospel-of-matthew}{composition}{catholic-traditional-v1}{disputed}{about the year 40-42}{c. A.D. 40--42}; \chronologyannotationreach{Matt.22.1}{Matt}{true}\chronologyannotationreach{Matt.22.10}{Matt}{true}\chronologyannotationreach{Matt.22.11}{Matt}{true}\chronologyannotationreach{Matt.22.12}{Matt}{true}\chronologyannotationreach{Matt.22.13}{Matt}{true}\chronologyannotationreach{Matt.22.14}{Matt}{true}\chronologyannotationreach{Matt.22.2}{Matt}{true}\chronologyannotationreach{Matt.22.3}{Matt}{true}\chronologyannotationreach{Matt.22.4}{Matt}{true}\chronologyannotationreach{Matt.22.5}{Matt}{true}\chronologyannotationreach{Matt.22.6}{Matt}{true}\chronologyannotationreach{Matt.22.7}{Matt}{true}\chronologyannotationreach{Matt.22.8}{Matt}{true}\chronologyannotationreach{Matt.22.9}{Matt}{true}\chronologyannotationclaim{composition.gospel-of-matthew}{composition}{catholic-traditional-v1}{disputed}{the years 40-45}{A.D. 40--45}; \chronologyannotationreach{Matt.22.1}{Matt}{true}\chronologyannotationreach{Matt.22.10}{Matt}{true}\chronologyannotationreach{Matt.22.11}{Matt}{true}\chronologyannotationreach{Matt.22.12}{Matt}{true}\chronologyannotationreach{Matt.22.13}{Matt}{true}\chronologyannotationreach{Matt.22.14}{Matt}{true}\chronologyannotationreach{Matt.22.2}{Matt}{true}\chronologyannotationreach{Matt.22.3}{Matt}{true}\chronologyannotationreach{Matt.22.4}{Matt}{true}\chronologyannotationreach{Matt.22.5}{Matt}{true}\chronologyannotationreach{Matt.22.6}{Matt}{true}\chronologyannotationreach{Matt.22.7}{Matt}{true}\chronologyannotationreach{Matt.22.8}{Matt}{true}\chronologyannotationreach{Matt.22.9}{Matt}{true}\chronologyannotationclaim{composition.gospel-of-matthew}{composition}{catholic-traditional-v1}{disputed}{about the year 60-68}{c. A.D. 60--68}; \chronologyannotationreach{Matt.22.1}{Matt}{true}\chronologyannotationreach{Matt.22.10}{Matt}{true}\chronologyannotationreach{Matt.22.11}{Matt}{true}\chronologyannotationreach{Matt.22.12}{Matt}{true}\chronologyannotationreach{Matt.22.13}{Matt}{true}\chronologyannotationreach{Matt.22.14}{Matt}{true}\chronologyannotationreach{Matt.22.2}{Matt}{true}\chronologyannotationreach{Matt.22.3}{Matt}{true}\chronologyannotationreach{Matt.22.4}{Matt}{true}\chronologyannotationreach{Matt.22.5}{Matt}{true}\chronologyannotationreach{Matt.22.6}{Matt}{true}\chronologyannotationreach{Matt.22.7}{Matt}{true}\chronologyannotationreach{Matt.22.8}{Matt}{true}\chronologyannotationreach{Matt.22.9}{Matt}{true}\chronologyannotationclaim{composition.gospel-of-matthew}{composition}{catholic-traditional-v1}{disputed}{about the years 64-67}{c. A.D. 64--67}; \chronologyannotationreach{Matt.22.1}{Matt}{true}\chronologyannotationreach{Matt.22.10}{Matt}{true}\chronologyannotationreach{Matt.22.11}{Matt}{true}\chronologyannotationreach{Matt.22.12}{Matt}{true}\chronologyannotationreach{Matt.22.13}{Matt}{true}\chronologyannotationreach{Matt.22.14}{Matt}{true}\chronologyannotationreach{Matt.22.2}{Matt}{true}\chronologyannotationreach{Matt.22.3}{Matt}{true}\chronologyannotationreach{Matt.22.4}{Matt}{true}\chronologyannotationreach{Matt.22.5}{Matt}{true}\chronologyannotationreach{Matt.22.6}{Matt}{true}\chronologyannotationreach{Matt.22.7}{Matt}{true}\chronologyannotationreach{Matt.22.8}{Matt}{true}\chronologyannotationreach{Matt.22.9}{Matt}{true}\chronologyannotationclaim{composition.gospel-of-matthew}{composition}{catholic-traditional-v1}{disputed}{about the year 50}{c. A.D. 50}.}%
  \space \chronologyannotationcomparisongroup{catholic-critical-v1}{composition}{composition-only}{\textbf{Comparison (catholic-critical-v1) -- Composition} -- composition-only: \chronologyannotationreach{Matt.22.1}{Matt}{true}\chronologyannotationreach{Matt.22.10}{Matt}{true}\chronologyannotationreach{Matt.22.11}{Matt}{true}\chronologyannotationreach{Matt.22.12}{Matt}{true}\chronologyannotationreach{Matt.22.13}{Matt}{true}\chronologyannotationreach{Matt.22.14}{Matt}{true}\chronologyannotationreach{Matt.22.2}{Matt}{true}\chronologyannotationreach{Matt.22.3}{Matt}{true}\chronologyannotationreach{Matt.22.4}{Matt}{true}\chronologyannotationreach{Matt.22.5}{Matt}{true}\chronologyannotationreach{Matt.22.6}{Matt}{true}\chronologyannotationreach{Matt.22.7}{Matt}{true}\chronologyannotationreach{Matt.22.8}{Matt}{true}\chronologyannotationreach{Matt.22.9}{Matt}{true}\chronologyannotationcomparisonclaim{critical.gospel-of-matthew}{composition}{catholic-critical-v1}{catholic-critical-v1}{preferred}{passage.united-states-conference-of-catholic-bishops.new-american-bible-revised-edition.english-usccb-web-2026-09-21.matthew-introduction}{post-A.D. 70 date}{The Greek Gospel of Matthew in the NABRE introduction, Post-A.D. 70 date}.}%
}
% offertory: Offertory
\expandafter\def\csname triptychchronologyannotation@offertory\endcsname{%
  \chronologyannotationgroup{composition}{preferred}{\textbf{Composition}: \chronologyannotationreach{Ps.137.7}{Ps}{true}\chronologyannotationclaim{critical.psalms.latest-composition-boundary}{composition}{catholic-critical-v1}{preferred}{before the Maccabean period, around 165 B.C.; no individual psalm can be dated securely}{Before c. 165 B.C.}}%
}
% communion: Communion
\expandafter\def\csname triptychchronologyannotation@communion\endcsname{%
  \chronologyannotationgroup{composition}{preferred}{\textbf{Composition}: \chronologyannotationreach{Ps.118.4}{Ps}{true}\chronologyannotationreach{Ps.118.5}{Ps}{true}\chronologyannotationclaim{critical.psalms.latest-composition-boundary}{composition}{catholic-critical-v1}{preferred}{before the Maccabean period, around 165 B.C.; no individual psalm can be dated securely}{Before c. 165 B.C.}}%
}


The appointed biblical units appear below in canonical and verse order.
The psalm dates are broad limits of composition, not dates of each
poem's particular occasion. The orations and the Introit's composed
antiphon have no biblical composition-date entry.

\begin{dossiertable}
Introit verse & Ps 77(78):1 & Israel; instruction addressed to the people & \chronodate{introit}{\chronologyannotation{introit}}\\
\dossierprose{The psalm teaches through Israel's deliverances and failures and ends with David's shepherding. Augustine hears God through the prophet; Bellarmine identifies David as immediate speaker. These are interpretations of the voice, not a dated occasion of composition. The USCCB introduction supplies the broad boundary and warns that individual psalms cannot be securely dated.}
Alleluia & Ps 104(105):1 & Israel's covenant memory; proclamation among nations & \chronodate{alleluia}{\chronologyannotation{alleluia}}\\
\dossierprose{The psalm recalls covenant, Joseph, Exodus, and land. Those recalled events are distinct from the writing of the psalm and its later Christian reception. No single composition site or individual life stage is established.}
Communion & Ps 118(119):4--5 & Prayer within Israel's life under God's law & \chronodate{communion}{\chronologyannotation{communion}}\\
\dossierprose{The speaker's wish for directed ways belongs to the opening of an alphabetic meditation on divine instruction. Its recurring prayers, afflictions, promises, and confessions do not identify one recoverable biographical occasion. The displayed boundary does not date every constituent of the Psalter together.}
Offertory & Ps 137(138):7 & Thanksgiving and trust amid hostile surroundings & \chronodate{offertory}{\chronologyannotation{offertory}}\\
\dossierprose{The psalm gives thanks for answered prayer and trusts God's continuing work. Neither the named trouble nor the enemies establish an individually dated event. The Offertory's verbal adaptation belongs to its liturgical use.}
Gradual & Ps 140(141):2 & Prayer imaged through incense and evening sacrifice & \chronodate{gradual}{\chronologyannotation{gradual}}\\
\dossierprose{The fuller prayer concerns speech, temptation, correction, and danger. Its sacrificial language supplies the image, not a securely located occasion of composition. Augustine's interpretation through Christ's Passion is distinct from the poem's composition chronology.}
Gospel & Matt 22:1--14 & Jesus teaching in the Temple, Jerusalem; traditional composition associated with Palestine & \chronodate{gospel}{\chronologyannotation{gospel}}\\
\dossierprose{Traditional attribution names Matthew. Jacquier reports several composition proposals; Durand's early proposal concerns a hypothesized Aramaic original and does not date the Greek rendering. The separately displayed critical comparison concerns the Greek Gospel: the USCCB introduction treats authorship and dependence on Mark as critical questions and suggests Antioch as a plausible location. The proposals concern distinguishable literary objects and arguments, not one agreed date.}
\dossierevent{Narrated setting: Matthew 21:23 places the controversy in the Temple; 21:45 identifies the chief priests and Pharisees, and 22:1 continues the speech. Jerusalem is the event's setting. No numerical date for this narrated occasion is supplied here.}
Epistle & Eph 4:23--28 & Traditional Pauline captivity; Rome preferred to Caesarea in Ladeuze's account; recipients associated with Ephesus in Asia Minor & \chronodate{epistle}{\chronologyannotation{epistle}}\\
\dossierprose{Ladeuze's captivity range and Prat's narrower Pauline chronology remain distinct alternatives. The USCCB Ephesians introduction, paragraphs 3--4, retains traditional Pauline captivity but reports doubts about direct authorship based on style, vocabulary, comparison with Colossians, and ecclesiology. It allows a secretary under Paul's direction or a later disciple developing his thought. The separate comparison's approximate interval belongs only to the later-disciple hypothesis; no composition city or separate secretarial date is supplied for those alternatives. The absence of personal greetings and omission of the Ephesian address in important early manuscripts support proposals for circulation among Asian churches or a Laodicean destination. These proposed recipients do not establish the writer's location.}
\end{dossiertable}

The Gospel's Temple encounter remains distinct from Matthew's composition
and transmission. The chronology witnesses listed in the references
supply the displayed dates and alternatives.
